% Part: second-order-logic
% Chapter: sol-and-sets
% Section: comparing-sets

\documentclass[../../../include/open-logic-section]{subfiles}

\begin{document}

\olfileid[hi]{sol}{set}{cmp}

\olsection{समुच्चयों की तुलना}

\begin{prop}
!!{formula} $\lforall[x][(X(x) \lif Y(x))]$ उपसमुच्चय-संबंध को परिभाषित
करता है, अर्थात् $\Sat{M}{\lforall[x][(X(x) \lif Y(x))]}[s]$ तभी और केवल
तभी जब $s(X) \subseteq s(Y)$।
\end{prop}

\begin{prop}
!!{formula} $\lforall[x][(X(x) \liff Y(x))]$ समुच्चयों पर तादात्म्य-संबंध
को परिभाषित करता है, अर्थात्
$\Sat{M}{\lforall[x][(X(x) \liff Y(x))]}[s]$ तभी और केवल तभी जब $s(X) =
s(Y)$।
\end{prop}

\begin{prop}
!!{formula} $\lexists[x][X(x)]$ अरिक्त होने के गुण को परिभाषित करता है,
अर्थात् $\Sat{M}{\lexists[x][X(x)]}[s]$ तभी और केवल तभी जब $s(X) \neq
\emptyset$।
\end{prop}

समुच्चय~$X$ का आकार समुच्चय~$Y$ से अधिक नहीं है, $\cardle{X}{Y}$, तभी और
केवल तभी जब कोई !!a{injective} फलन $f\colon X \to Y$ हो। चूँकि हम व्यक्त
कर सकते हैं कि कोई फलन अंतःक्षेपी है, और यह भी कि~$X$ के अवयवों पर उसके
मान~$Y$ में हैं, इसलिए हम प्रांत के उपसमुच्चयों पर ``आकार में अधिक न होने''
के संबंध को भी परिभाषित कर सकते हैं।

\begin{prop}
यह सूत्र
\[
\lexists[u][(\lforall[x][(X(x) \lif Y(u(x)))] \land \lforall[x][\lforall[y][(\eq[u(x)][u(y)] \lif
      \eq[x][y])]])]
\]
आकार में अधिक न होने के संबंध को परिभाषित करता है।
\end{prop}

दो समुच्चय समान आकार के, अथवा ``समसंख्यक'', $\cardeq{X}{Y}$, तभी और केवल
तभी होते हैं जब कोई !!a{bijective} फलन~$f\colon X \to Y$ हो।

\begin{prop}
यह सूत्र
\begin{multline*}
  \lexists[u][(\lforall[x][(X(x) \lif Y(u(x)))] \land {}\\
    \lforall[x][\lforall[y][(\eq[u(x)][u(y)] \lif \eq[x][y])]]
  \land {} \\\lforall[y][(Y(y) \lif \lexists[x][(X(x)
      \land \eq[y][u(x)])])])]
\end{multline*}
समसंख्यकता के संबंध को परिभाषित करता है।
\end{prop}

हम इन !!{formula} को क्रमशः $X \subseteq Y$, $X = Y$, $X \neq
\emptyset$, $\cardle{X}{Y}$ और $\cardeq{X}{Y}$ से संक्षिप्त करेंगे। (यह
कुछ भ्रम उत्पन्न कर सकता है, क्योंकि समुच्चयों $X$ और $Y$ की अनौपचारिक
चर्चा में भी हम यही संकेत-लिपि उपयोग करते हैं---पर यहाँ यह संकेत-लिपि
एक-स्थानी संबंध-चरों $X$ और~$Y$ वाले द्वितीय-कोटि तर्कशास्त्र के !!{formula}
का संक्षिप्त रूप है।)

\begin{prop}
!!{sentence} $\lforall[X][\lforall[Y][((\cardle{X}{Y} \land
    \cardle{Y}{X}) \lif \cardeq{X}{Y})]]$ वैध है।
\end{prop}

\begin{proof}
यह !!{sentence} किसी !!a{structure}~$\Struct{M}$ में तब संतुष्ट होता है
जब किन्हीं उपसमुच्चयों $X \subseteq \Domain{M}$ और $Y \subseteq
\Domain{M}$ के लिए, यदि $\cardle{X}{Y}$ और $\cardle{Y}{X}$, तो
$\cardeq{X}{Y}$। किंतु यह \emph{किन्हीं भी} समुच्चयों $X$ और $Y$ के लिए
लागू होता है---यही श्रेडर--बर्नस्टाइन प्रमेय है।
\end{proof}


\end{document}
